\documentclass[12pt]{article}
\usepackage{amsmath}
\usepackage{amsfonts}
\usepackage{amssymb}
\usepackage{physics}
\usepackage{geometry}
\usepackage{parskip}
\geometry{a4paper, margin=0.7in}


\title{02YMECH - Solution 7}
\date{Assigned for the week of Nov 3, 2025}
\author{}

\begin{document}
\maketitle

\section*{Solutions}

\[
r(\theta)=\frac{a}{1+b\cos\theta}, \qquad a>0,\;0<b<1.
\]

As we derived in lectures:

\[
f(r)=-\frac{l^2u^2}{m} \!\left(\frac{d^2u}{d\theta^2}+u\right),
\]
where \(u=1/r\). Then
\[\
u=\frac{1}{a}(1+b\cos\theta),\qquad
\frac{d^2u}{d\theta^2}+u=\frac{1}{a}.
\]


\begin{enumerate}
\item \textbf{Central force}
\[
f(r)=-\frac{l^2}{ma}\,\frac{1}{r^2}.
\]

Thus the force is an inverse--square law.

\item \textbf{Attractive or repulsive}

Since \(m>0\), \(l^2>0\), and \(a>0\),
\[
f(r)<0 \quad \text{for all } r>0,
\]
so the force is \emph{attractive}, in particular for \(0<r<a\).
\end{enumerate}

\end{document}
